\section{Conclusion}
\label{sec:conclusion}

The bid-rigging detection literature has been built on the
assumption that bid-level data are observable. Real enforcement
environments routinely violate that assumption: many electronic
procurement systems preserve only the contract-award envelope,
and the bid-distribution methods that work on the bid layer
become inoperative on the layer that survives. The paper takes
this informational mismatch as its starting point and asks what
collusion-relevant statistic can be recovered when only the
award layer is observable.

Endogenous loser-side participation is such a statistic. A
separating-equilibrium framework with cover bidders and genuine
entrants identifies $\log(1+\text{tenders\_count})$ together
with the $\text{wins}=0$ filter as a sufficient ranking
statistic for cover-bidder type given award-record data; the
binary $\text{median}+1.5\!\times\!\text{IQR}$ rule we deploy
is its information-coarsening. On São Paulo's BEC
($\valYearStart$--$\valYearEnd$), the construct discriminates
$\valCobidders$ adjudicated CADE cobidders inside the always-loser
stratum at firm-level AUC $\valAUCFLfirmTemp$ under temporal
holdout (train $2009$--$2016$, test $2017$--$2019$) and
$\valAUCprePost$ on a strict pre-$2020$ contemporaneous benchmark
with participation-stratified permutation null
($\valMultThreeTwo\!\!\times$ excess over the random-matching
baseline, $p<0.001$). Against bid-distribution detectors that
require per-bidder bid amounts the screening statistic reaches
comparable accuracy on a substantially thinner data envelope (AUC
$\valImhofFLBin$ vs.\ $\valImhofFull$) and adds non-redundant
signal in same-sample combination ($+\valImhofIncFL$ AUC, DeLong
$p=\valImhofIncFLp$). Within the always-loser stratum, the
cobidder validation positive class matches the modeled cover-
bidder type along four of five operational predictions, with
Cohen's $d$ up to $+1.00$. AUC against the $\valDirectCADE$
direct CADE defendants in the broader BEC firm universe is
$\valAUCdirectStd$; the cobidder/direct-defendant asymmetry is
the design's empirical signature, since the construct recovers
loser-side participation rather than winner-side identity.

The pricing evidence is descriptive corroboration. Frequent-loser
presence associates with $\valHeadlineRange$ higher conditional
log negotiated unit price across four estimators on the broad
sample. The overlap-restricted estimate is $\valMatchOverlapCoef$
under ATT-reweighting and $\valBetaUnwOverlap$ under unweighted
overlap restriction; the empirical decomposition in
\S\ref{sec:results_overlap} attributes the sign reversal to
ATT-weights up-weighting cells with relatively few untreated
counterfactuals rather than to the overlap restriction itself
($\valDroppedItems$ of $79{,}452$ treated items, or
$\valDroppedItemsShare$, are strictly dropped). We report all
specifications and acknowledge the screening-value reading is one
interpretation among several the data do not adjudicate. The
discrimination evidence and the architectural test do not depend
on overlap weighting, and they carry the contribution.

Three implications for enforcement design follow.

\paragraph{What survives data coarsening} The information that
the bid-layer literature reads off the bid distribution is
partially recoverable from a different statistic at the award
layer when bid data are unavailable. The screening object the
recovery delivers is a participation statistic, not a price
moment, and the equilibrium-generated participation footprint
that the collusive arrangement leaves on the award layer is the
appropriate empirical object under coarsened observability.

\paragraph{Screening and forensic stages are sequenced, not
parallel} The architectural test in §\ref{sec:forensic} reads
the screening statistic and the bid-distribution moment battery
as informational complements: comparable accuracy on different
data envelopes (AUC $\valImhofFLBin$ vs $\valImhofFull$ on the
same cartel-adjacency target) and non-redundant signal when
both envelopes are available (combined AUC
$\valImhofComboCont$, $\valImhofIncCombo$ over either alone). Under incomplete observability the deployable
architecture is therefore an award-layer screening stage that
produces candidate environments, followed by a bid-layer
forensic stage that runs the bid-distribution moment battery on
those candidates where microdata are recoverable. The two
stages answer different questions and require different
statistics.

\paragraph{Proactive screening does not require bid-microdata
archival as a precondition} Reforms that mandate bid archival
expand the forensic stage; they do not substitute for the
screening one. Jurisdictions that lack the bid layer can deploy
the award-layer stage on data they already maintain, and the
sequencing reading reverses the standard policy intuition that
bid-archival is the rate-limiting precondition for proactive
cartel screening
\citep{becker1968crime,stigler1970optimum,baker2002case,harrington2008detecting}.

\paragraph{Scope and an agenda}
The contribution stops at the screening stage by design.
Award-record data does not contain what membership identification
requires, and the cobidder population is the validation object
the data layer supports. Adversarial adaptation degrades the construct's
discrimination predictably (Online Appendix~G), and replication
rests on three diagnostic conditions---a winner/loser
participation registry, an electronic platform, and an
institutional asymmetry between the two sides of the bidding
pool---all observable \emph{ex ante} in any candidate
jurisdiction. Three open questions follow naturally.
\emph{First,} cross-jurisdictional replication in
ComprasNet, EU TED, and U.S.~contracting-officer files,
particularly in pregão-style real-time-electronic auctions
where the modal contrast in §\ref{sec:mech_modal} suggests the
screening signal is sharpest. \emph{Second,} extension beyond
commodity-heavy item composition to services, infrastructure,
and public works, where within-product-code price comparison
is harder but the participation primitive is unchanged.
\emph{Third,} formal characterization of the optimal sequencing
of screening and forensic stages under endogenous adversarial
adaptation, building on the framework's deployment problem.

The conceptual content travels: wherever an enforcement
environment exposes the award layer without exposing the bid
layer, the question is not whether bid-distribution screens can
be retrofitted onto data that does not contain them, but what
participation primitive the collusive arrangement generates that
survives the coarsening, and how the screening stage that reads
that primitive should be sequenced before the forensic stage
that requires bid-level evidence.
